% NeurIPS 2026 Paper Checklist
% Mandatory — missing checklist results in desk rejection.

\section*{NeurIPS Paper Checklist}

\begin{enumerate}

\item {\bf Claims}
    \begin{enumerate}
        \item Do the main claims made in the abstract and introduction
              accurately reflect the paper's contributions and scope?
    \answerYes{The abstract states two specific research questions
    (cross-generational transfer and safety scaling) and the introduction
    enumerates four concrete contributions, each tied to a results section.}
    \end{enumerate}

\item {\bf Limitations}
    \begin{enumerate}
        \item Does the paper discuss the limitations of the work performed
              by the authors?
    \answerYes{Section~\ref{sec:discussion} includes a dedicated Limitations
    paragraph covering: small N for scaling, single model family,
    HarmBench coverage bounds, LLM judge bias, and evolutionary probe
    coverage limitations.}
    \end{enumerate}

\item {\bf Theory Assumptions and Proofs}
    \begin{enumerate}
        \item For each theoretical result, does the paper provide the full
              set of assumptions and a complete (or correct) proof?
    \answerNA{This is an empirical paper with no theoretical results.}
    \end{enumerate}

\item {\bf Experimental Result Reproducibility}
    \begin{enumerate}
        \item Does the paper fully disclose all the information needed to
              reproduce the main experimental results of the paper to the
              extent that it affects the main claims and/or conclusions of
              the paper (regardless of whether the code and data are
              provided or not)?
    \answerYes{Section~\ref{sec:setup} specifies all models, hyperparameters
    (grid dimensions, batch size, generations, concurrency), seed dataset
    (HarmBench 400 behaviors), judge model (Llama-3.1-8B-Instruct,
    self-hosted), serving
    infrastructure (vLLM 0.19.1, bf16, temperature 0), and hardware
    (H200, RTX~4090). The
    transfer replay protocol is fully described in
    Section~\ref{sec:method:transfer}.}
    \end{enumerate}

\item {\bf Open access to data and code}
    \begin{enumerate}
        \item Does the paper provide open access to the data and code, with
              sufficient instructions to faithfully reproduce the main
              experimental results, as described in supplemental material?
    \answerYes{Anonymous code repository is provided as supplemental
    material, including the full evolution framework (Rust), analysis
    scripts (Python), orchestration scripts, and instructions for
    reproducing all experiments. Seed behaviors are from the public
    HarmBench dataset. Experiment result JSONs are included.}
    \end{enumerate}

\item {\bf Experimental Setting/Details}
    \begin{enumerate}
        \item Does the paper specify all the training and test details
              (e.g., data splits, hyperparameters, how they were chosen,
              type of optimizer, etc.) necessary to understand the results?
    \answerYes{All MAP-Elites hyperparameters are specified in
    Section~\ref{sec:setup}. No training is performed; we evaluate
    existing models. Evolution configuration: 200 generations,
    batch size 50, $6 \times 6 \times 8$ MAP-Elites grid, mutation
    rate 0.3, crossover rate 0.7.}
    \end{enumerate}

\item {\bf Experiment Statistical Significance}
    \begin{enumerate}
        \item Does the paper report error bars suitably and correctly
              defined or other appropriate information about the statistical
              significance of the experiments?
    \answerYes{All experiments use 3 independent random seeds per model.
    We report mean $\pm$ standard deviation and 95\% bootstrap confidence
    intervals. Paired bootstrap tests are used for pairwise model
    comparisons ($p < 0.05$ threshold). We acknowledge that $N = 3$
    provides limited power for some contrasts.}
    \end{enumerate}

\item {\bf Experiments Compute Resources}
    \begin{enumerate}
        \item For each experiment, does the paper provide sufficient
              information on the computer resources (type of compute
              workers, memory, time for each experiment, etc.) needed to
              reproduce the experiments?
    \answerYes{Section~\ref{sec:setup} reports: NVIDIA H200 (141\,GB) for
    targets, RTX~4090 (24\,GB) for the 8B judge, vLLM 0.19.1. Total GPU
    time: $\sim$60 hours (12 evolution runs at $\sim$4h each, 36 transfer
    replays at $\sim$8min each, 4 raw-seed baselines, 1 judge audit with
    Llama-3.1-70B on the H200).}
    \end{enumerate}

\item {\bf Code Of Ethics}
    \begin{enumerate}
        \item Does the research conducted in the paper conform, in every
              respect, with the NeurIPS Code of Ethics?
    \answerYes{This work studies LLM safety alignment to improve defenses.
    We follow responsible disclosure practices: we do not release
    specific jailbreak prompts that could enable misuse, only aggregate
    statistics and methodology. The goal is to help model developers
    identify and fix vulnerabilities.}
    \end{enumerate}

\item {\bf Broader Impacts}
    \begin{enumerate}
        \item Does the paper discuss both potential positive societal
              impacts and negative societal impacts of the work performed?
    \answerYes{Positive: enables systematic safety evaluation across model
    generations, helping developers prioritize alignment work. Negative:
    methodology could in principle be used to find and exploit
    vulnerabilities. We mitigate this by not publishing specific attack
    prompts and by framing the work as a defensive evaluation tool.}
    \end{enumerate}

\item {\bf Safeguards}
    \begin{enumerate}
        \item Does the paper describe safeguards that have been put in
              place for responsible release of data or models that have a
              high risk for misuse (e.g., pretrained language models,
              image generators, or scraped datasets)?
    \answerYes{We release aggregate statistics (transfer rates, ASR per
    category) and the evolution framework, but not the specific attack
    prompts or model responses from our experiments. The seed behaviors
    come from the publicly available HarmBench dataset.}
    \end{enumerate}

\item {\bf Licenses for existing assets}
    \begin{enumerate}
        \item Are the creators or original owners of assets (e.g., code,
              data, models), used in the paper, properly credited and are
              the license and terms of use explicitly mentioned and properly
              respected?
    \answerYes{HarmBench is cited and used under its original license.
    Gemma models are used under Google's Gemma license. The judge model
    (Llama-3.1-8B-Instruct) is used under Meta's license.}
    \end{enumerate}

\item {\bf New Assets}
    \begin{enumerate}
        \item Are new assets introduced in the paper well documented and
              is the documentation provided alongside the assets?
    \answerYes{The evolution framework is documented with
    README, CLI help text, and reproduction instructions. Analysis
    scripts include docstrings and usage examples.}
    \end{enumerate}

\item {\bf Crowdsourcing and Research with Human Subjects}
    \begin{enumerate}
        \item For crowdsourcing experiments and research with human
              subjects, does the paper include the full text of instructions
              given to participants and screenshots, if applicable, as well
              as details about compensation?
    \answerNA{No human subjects or crowdsourcing.}
    \end{enumerate}

\item {\bf Institutional Review Board (IRB) Approvals or Equivalent for
         Research with Human Subjects}
    \begin{enumerate}
        \item Does the paper describe potential risks incurred by study
              participants, whether combative or non-combative to the
              subjects, and does it meet the necessary IRB criteria?
    \answerNA{No human subjects.}
    \end{enumerate}

\item {\bf Declaration of LLM usage}
    \begin{enumerate}
        \item Does the paper describe the usage of LLMs in the research
              and/or writing process?
    \answerYes{LLMs are used in two capacities: (1)~as evaluation targets
    (the Gemma models being probed for safety), and (2)~as a self-hosted
    automated judge (Llama-3.1-8B-Instruct scoring target responses).}
    \end{enumerate}

\end{enumerate}
